% Ch2_4, slide 4: moving a point charge q by dl in the field E: the force F_E = qE, its component
% F_EL along a_l
\begin{tikzpicture}[scale=1.35]
  \foreach \x in {-1.6,-0.8,0,0.8,1.6}{\draw[ELcField!55,line width=0.7pt,-{Stealth[length=5pt]}] (\x-1.0,-1.6)--(\x+0.9,1.7);}
  \node[ELlab,text=ELcField,anchor=north] at (-2.3,-1.65) {$\vect{E}$};
  \draw[ELdash] (-3.0,0)--(2.6,0);
  \coordinate (q) at (-0.05,0);
  \draw[ELvec2,line width=1.2pt] (q)--++(0.62,1.1) node[ELlabs,anchor=west]{$\vect{F}_E$};
  \draw[ELvec3,line width=1.2pt] (q)--++(-0.62,0) node[ELlabs,anchor=south]{$\uvec{\ell}$};
  \draw[ELdash,ELcThird] (0.57,1.1)--(0.57,0);
  \draw[ELvec3,line width=1.2pt] (q)--++(0.62,0) node[ELlabs,anchor=north west]{$F_{EL}$};
  \fill[ELink] (q) circle (2.2pt); \node[ELlabs,anchor=north] at (0,-0.08) {$q$};
  \draw[ELthin,{Stealth[length=4pt]}-{Stealth[length=4pt]}] (-1.0,-0.45)--(-0.1,-0.45) node[ELlabs,anchor=north,pos=0.5]{$d\ell$};
\end{tikzpicture}
